\documentclass[11pt]{article}
\usepackage[utf8]{inputenc}
\usepackage{amsmath,amssymb}
\usepackage[margin=1in]{geometry}
\usepackage{hyperref}
\emergencystretch=3em

\title{Normalised remainders for tight families of Taylor data}
\author{Claude, with Daniel Murfet}
\date{2026-09-07}

\begin{document}
\maketitle
\sloppy

\begin{abstract}
The localisation upgrade suggested by Astra~\#11. The bounded-family theorems of unit~150 required a
deterministic bound $\|X_n\|\le M$ on the random Taylor data, which excludes Gaussian-type limits.
Here the hypothesis is replaced by tightness of the norms,
$\forall\eta>0\ \exists M\ \forall^\infty n,\ \mu\{\|X_n\|>M\}\le\eta$ (the content of the paper's
appendix prop:convergence~(i)). On the ball the error is deterministically small; off the ball the
probability is at most $\eta$; so the error tends to zero in probability
(\texttt{tendstoInMeasure\_zero\_of\_tight}), and Mathlib's
\texttt{tendstoInDistribution\_of\_tendstoInMeasure\_sub} yields
$\mathrm{normA}\Rightarrow A_\alpha(Z)$ and $\mathrm{normB}\Rightarrow B_\alpha(Z)$
(\texttt{tendstoInDistribution\_normA\_tight}, \texttt{tendstoInDistribution\_normB\_tight}).
Bounded families are tight (\texttt{tight\_of\_bounded}), so unit~150 is a special case. Zero
\texttt{sorry}/\texttt{axiom}.
\end{abstract}

\section{The things formalised}
{\small
\begin{verbatim}
theorem tendstoInMeasure_zero_of_tight (μ) (f : ι → Ω → ℝ) (X : ι → Ω → CoeffPair)
    (e : ℝ → ι → ℝ)
    (he : ∀ M, 0 ≤ M → Tendsto (e M) l (𝓝 0))
    (hf : ∀ M, 0 ≤ M → ∀ᶠ n in l, ∀ ω, ‖X n ω‖ ≤ M → |f n ω| ≤ e M n)
    (htight : ∀ η : ENNReal, 0 < η → ∃ M, 0 ≤ M ∧ ∀ᶠ n in l, μ {ω | M < ‖X n ω‖} ≤ η) :
    TendstoInMeasure μ f l (fun _ => 0)
theorem tendstoInDistribution_normA_tight ... (X) (hXm : ∀ n, Measurable (X n))
    (htight : ∀ η : ENNReal, 0 < η → ∃ M, 0 ≤ M ∧ ∀ᶠ n in l, μ {ω | M < ‖X n ω‖} ≤ η)
    (Z) (hX : TendstoInDistribution X l Z (fun _ => μ) μ') (Nseq)
    (hN : Tendsto Nseq l atTop)
    (α) (hα : α ∈ polesBelowGen h₁ h₂ k₁ k₂ p₁ p₂ T) :
    TendstoInDistribution (fun n ω => normA … α (Nseq n) (X n ω)) l
      (coeffA … α ∘ Z) (fun _ => μ) μ'
theorem tendstoInDistribution_normB_tight ... : … normB … ⇒ coeffB … α ∘ Z
theorem tight_of_bounded (μ) (X) (M) (hM : 0 ≤ M) (hXb : ∀ n ω, ‖X n ω‖ ≤ M) : (tightness)
\end{verbatim}
}

\section{Mathematics}
Fix $\varepsilon>0$ and $\eta>0$; choose $M$ from tightness. For $n$ large,
$\mathrm{errA}_M(N_n)<\varepsilon$, so $\{|\mathrm{err}_n|\ge\varepsilon\}\subseteq\{\|X_n\|>M\}$
and $\mu\{|\mathrm{err}_n|\ge\varepsilon\}\le\eta$. Since $\eta$ is arbitrary, $\mathrm{err}_n\to0$
in probability. The rest is the same as unit~150, with Mathlib's
\texttt{tendstoInDistribution\_of\_tendstoInMeasure\_sub} replacing the additive Slutsky lemma.

\section{Paper-facing meaning}
This is the form of the ordered normalised-remainder convergence in thm:strataempiricalexpansion
that applies to the empirical process: the paper establishes tightness of $\tilde\xi_n$ from the
moment bound $\mathbb E\sup|\tilde\xi_n|^s\le C$ (appendix lem:thm58new and prop:convergence~(i));
in our formulation the corresponding hypothesis is norm tightness of the Taylor data in
\texttt{CoeffPair}, which follows from the paper's bound by Cauchy estimates (not formalised) or
can be taken as the standing hypothesis. Units~145--152 together give, for $d=2$ and one chart:
coefficient continuity, convergence in distribution of finitely many coefficients, convergence of
the ordered normalised remainders for tight families, and exponential negligibility of the far
region.

\section{Lean notes}
\begin{itemize}
\item \texttt{ENNReal.tendsto\_nhds\_zero}: $u\to0$ iff $\forall\eta>0,\ \forall^\infty n,\ u_n\le\eta$,
which is exactly the shape a tightness argument produces; \texttt{measure\_mono} on the set
inclusion finishes.
\item \texttt{tendstoInDistribution\_of\_tendstoInMeasure\_sub} takes the difference as the Pi
subtraction \texttt{Y - X}; a term stated with an explicit lambda difference is accepted by
\texttt{exact} (definitional unfolding of \texttt{Pi.sub\_apply}).
\item \texttt{zero\_le} for \texttt{ENNReal} needs a type ascription
(\texttt{(zero\_le : (0 : ENNReal) ≤ η)}); \texttt{zero\_le \_} fails since the argument is implicit.
\end{itemize}

\end{document}
